\documentclass[10pt,a4paper]{amsart}
\usepackage[margin=19mm]{geometry}
\usepackage{amsmath,amssymb,mathtools}
\usepackage[scr=rsfs,cal=euler]{mathalfa}
\usepackage{xfrac}
\usepackage[unicode,hidelinks,bookmarks=false]{hyperref}
\newcommand{\ie}{i.e.\ }
\newcommand{\cf}{cf.\ }
\newcommand{\resp}{respectively\ }
\newcommand{\Subsection}{\subsection}
\providecommand{\nobreakdash}{\nobreak\hbox{-}}
\setlength{\emergencystretch}{2em}
\multlinegap0pt
\usepackage{amssymb}
\usepackage{enumerate}
\usepackage{xcolor}
\usepackage{tikz}
\newtheorem{lemma}{Lemma}
\def\ol{\overline}
\title{An introduction to separated graphs and their type semigroups}
\author{Pere Ara}
\date{Source: arXiv:2604.18304v1 (CC BY 4.0)}
\begin{document}
\maketitle

\noindent\emph{Source and licence.} An introduction to separated graphs and their type semigroups. Version arXiv:2604.18304v1 (CC BY 4.0), distributed by arXiv under the Creative Commons Attribution 4.0 International licence, http://creativecommons.org/licenses/by/4.0/, as recorded in the arXiv licence metadata for this exact version and shown on the abstract page https://arxiv.org/abs/2604.18304v1. Full licence notice: \texttt{licenses/arxiv-2604.18304-CC-BY-4.0.txt}.\par\smallskip

\noindent\textit{Editorial scope.} Counted unit: the article's lemma lem:coinvariants (source0.tex lines 964--968). The article's complete original proof of it is delivered with it (source0.tex lines 971--978, one \texttt{proof} environment of 8 lines). No mathematics is written, repaired or re-derived: the statement and the proof are the article's own lines, in the article's own notation, and no retained line is substituted or re-flowed.  No further source-local context is needed: every cross-reference of every retained line resolves inside the two delivered spans.  Nothing was omitted from any retained span, so the package carries no omission record.  No retained line contains a citation, so the wrapper carries no bibliography and no bibliography entry.  No retained line contains a figure environment, an image-inclusion command or drawing code, so no image is delivered and the package declares no asset dependency.  Editorial formatting only, in addition: an AMS \texttt{amsart} preamble that restates the article's own declarations and macros used by the retained lines, namely the environments \texttt{lemma} on separate counters, together with the standard packages those lines need.  The retained environments print in this document as Lemma 1 in order of appearance, whereas the article prints the same items with the numbering of its own full document.  The complete native source of the article as distributed in the arXiv e-print for this exact version is bundled unchanged as \texttt{sources/198787/source0.tex}.
\begin{lemma}
	\label{lem:coinvariants}
	Let $\Omega$ be a complete set of representatives for the $G$-action on $E^0$.
	Consider the graph $E_G$ such that $E_G^0=E^0/G$ and $E_G^1 = \bigsqcup_{v\in \Omega} r^{-1} (v)$, with $r_G(e)=[r(e)]$ and $s_G(e) = [s(e)]$ for all $e\in E_G^1$. Then $M(E)_G\cong M(E_G)$ and therefore $M(E)_G$ is a graph monoid.
\end{lemma}

\begin{proof}  
	Define $\phi\colon M(E_G)\to M(E)_G$ by $\phi (a_{[v]}) = [a_v]\in M(E)_G$. The relations of $M(E_G)$ are clearly preserved. Indeed, if $v\in E^0$ and $\ol{v}\in \Omega$ is the unique element of $\Omega$ such that $[v]=[\ol{v}]$ in $E^0/G$, then we have
	$$\phi (a_{[v]}) = [a_v] = [a_{\ol{v}}] = [\sum _{e\in r^{-1}(\ol{v})} a_{s(e)}] = \sum_{e\in r_G^{-1}([v])} [a_{s(e)}] = \sum _{e\in r_G^{-1}([v])} \phi (s_G (e)).$$
	Hence $\phi$ defines a monoid homomorphism. Conversely, define $\psi \colon M(E)\to M(E_G)$
	by $\psi (a_v)= a_{[v]}$. The relations of $M(E)$ are preserved because every  $g\in G$ induces a bijection between $r^{-1}(v)$ and $r^{-1}(gv)$, for all $v\in E^0$. To see that $\psi$ factors through $M(E)_G$, it suffices to check that $\psi (a_v) = \psi (g\cdot a_v)$. But this is obvious, because $g\cdot a_v = a_{gv}$. 
    Hence we obtain a homomorphism $\ol{\psi}\colon M(E)_G\to M(E_G)$. 
    The maps $\phi$ and $\ol{\psi}$ are clearly mutually inverse.   
\end{proof}

\end{document}
